\documentclass[track=junior,level=3]{kidscode}
%% a picture to circle: \pick{icon}{word}
\newcommand{\pick}[2]{\begin{tabular}{@{}c@{}}{\fontsize{34}{40}\selectfont #1}\\[4pt]\textbf{#2}\end{tabular}}
%% a row of faces to circle for the child self-rating
\newcommand{\faces}{{\LARGE\faFrown[regular]\quad\faMeh[regular]\quad\faSmile[regular]}}
\newenvironment{howami}{\par\noindent\renewcommand{\arraystretch}{1.6}%
  \tabularx{\linewidth}{@{}>{\RaggedRight}X>{\centering\arraybackslash}p{4.6cm}@{}}
  \toprule \rowcolor{kcaccent!15}\textbf{Read it aloud} & \textbf{Circle a face}\\ \midrule}
  {\bottomrule\endtabularx\par}
\newcommand{\howrow}[1]{#1 & \faces\\}
\begin{document}
\kctitle{Show-Me Checklists}{Week 1 warm-up, week 5 and week 10 show-me checklists. One copy per child}

\begin{grownup}[These are not tests]
Nothing here is a test or a grade. You watch your child play and build, and check off what you
see. Print one copy of each page per child, in the week you use it. If a row did not come up
naturally, ask \enquote{Can you show me how you\dots?} at the end of the session. The Answer Key
(\enquote{How to use the show-me checklists}) says what each column looks like and what to do next.
\end{grownup}

\begin{center}\renewcommand{\arraystretch}{1.4}
\begin{tabularx}{\linewidth}{@{}>{\bfseries}p{3.4cm}X@{}}
\toprule
Not yet & Needs you to show it, or does not try it yet. That is fine: it is new.\\
Getting there & Does it with a hint, a reminder or a look at the book.\\
Got it & Does it alone, and can show it again another day.\\
\bottomrule
\end{tabularx}
\end{center}

%% =====================================================================
\kcpart{Week 1}{Show Me: Welcome Back}{A 5-minute warm-up on the Mac. Not graded. What carried over from Level 2?}

Nickname: \rule{6cm}{0.4pt}\hfill Date: \rule{3.5cm}{0.4pt}

\begin{grownup}[How to run the warm-up]
Open Scratch with a new project. Ask for one thing at a time: \enquote{Can you show me how you
make the cat\dots?} Two or three rows per child are plenty; skip the rest. Nobody fails: you only
want to know what to remind them of in weeks 1 and 2. With 3 or 4 children, let them show in pairs.
\end{grownup}

\begin{showme}[Week 1 warm-up]
\showmerow{keep a sprite moving with a forever loop}
\showmerow{make a sprite react with \enquote{if touching}}
\showmerow{keep score with a variable: set it to 0, change it by 1}
\showmerow{use a My Block or a message}
\showmerow{draw a square with the pen and a repeat}
\showmerow{find and fix a bug}
\end{showme}

\section{What would you like to make? Circle all you like.}
\begin{center}
%% alt: Five pictures to circle in a row: a game controller for games, a paintbrush for art, a bar chart for graphs, a robot for robots that learn and a code sign for writing code.
\begin{tikzpicture}
  \foreach \icon/\word [count=\i] in {\faGamepad/games,\faPaintBrush/art,\faChartBar/graphs,\faRobot/robots,\faCode/code}
    \node[align=center] at (\i*3.3,0) {\pick{\icon}{\word}};
\end{tikzpicture}
\end{center}

\section{How do you feel about writing code with words? Circle one.}
\begin{center}
\renewcommand{\arraystretch}{1.6}
\begin{tabular}{cccc}
{\fontsize{34}{40}\selectfont\faFrown[regular]} & {\fontsize{34}{40}\selectfont\faMeh[regular]} & {\fontsize{34}{40}\selectfont\faSmile[regular]} & {\fontsize{34}{40}\selectfont\faGrinStars[regular]}\\
worried & not sure & happy & super excited\\
\end{tabular}
\end{center}

\textbf{Notes for next time} (for the grown-up):
\answerspace{3.5cm}

%% =====================================================================
\kcpart{Week 5}{Show Me: Play Day}{Watch during the play day. Check one box in each row.}

Nickname: \rule{6cm}{0.4pt}\hfill Date: \rule{3.5cm}{0.4pt}
\medskip

\begin{showme}[Week 5]
\showmerow{use a repeat inside a repeat to make a pattern}
\showmerow{plan a program in parts before building it}
\showmerow{tally data and make a bar graph}
\showmerow{explain that a computer learns from examples and can learn unfair mistakes}
\showmerow{match Scratch blocks to Python words}
\showmerow{find and fix a bug}
\end{showme}

\medskip
\textbf{Notes for next time} (what to practice, what to celebrate):
\answerspace{4cm}

\section{How am I doing?}
\begin{grownup}
Read each line aloud. Your child circles a face. Every face is a good answer.
\end{grownup}
\begin{howami}
\howrow{I can make a pattern with a repeat inside a repeat.}
\howrow{I can plan a program in parts.}
\howrow{I can make a bar graph.}
\howrow{I can say how a computer learns from examples.}
\howrow{I can find the Python twin of a Scratch block.}
\howrow{I can fix a bug.}
\end{howami}

%% =====================================================================
\kcpart{Week 10}{Show Me: Graduation}{Watch during project time and the showcase. Check one box in each row.}

Nickname: \rule{6cm}{0.4pt}\hfill Date: \rule{3.5cm}{0.4pt}
\medskip

\begin{showme}[Week 10]
\showmerow{run a Python program in Thonny}
\showmerow{change numbers and colors in Python to make it their own}
\showmerow{use a \texttt{for} loop with an indent}
\showmerow{fix a typo error in Python}
\showmerow{plan on a plan card and build it}
\showmerow{explain one part of their code}
\showmerow{keep personal information private and give credit}
\end{showme}

\medskip
\textbf{Notes for next time} (to pass on to main-track Level 1):
\answerspace{4cm}

\section{How am I doing?}
\begin{grownup}
Read each line aloud. Your child circles a face. Look back at week 5 together: what can you do now?
\end{grownup}
\begin{howami}
\howrow{I can press Run in Thonny.}
\howrow{I can change numbers and colors in Python.}
\howrow{I can make a loop with \texttt{for}.}
\howrow{I can fix a typo.}
\howrow{I can plan on a plan card.}
\howrow{I can tell someone what my code does.}
\howrow{I keep my name and address private.}
\end{howami}

\end{document}
